\section*{अध्याय \ref{ch:b2:structure-determination} --- संरचना निर्धारण: \texorpdfstring{$^{13}$C}{13C} NMR, MS और IR एक साथ}

\begin{solution}{exo:b2:structure-determination:1}
1,2-डाइमेथिलबेन्ज़ीन: 4 (\ce{CH3}; C1/C2; C3/C6; C4/C5)। 1,3-: 5 (\ce{CH3}; C1/C3; C2; C4/C6; C5)।
1,4-: 3 (\ce{CH3}; C1/C4; C2/C3/C5/C6)।
\end{solution}

\begin{solution}{exo:b2:structure-determination:2}
\ce{CH3CH(OH)CH2CH3}। DEPT-135: C1 (\ce{CH3}) ऊपर, C2 (CH) ऊपर, C3 (\ce{CH2}) नीचे, C4 (\ce{CH3}) ऊपर।
DEPT-90: केवल C2।
\end{solution}

\begin{solution}{exo:b2:structure-determination:3}
209: कीटोन कार्बोनिल; 170: एस्टर (या अम्ल) कार्बोनिल; 129: \omterm{def:b2:huckel:aromatic}{ऐरोमैटिक} कार्बन; 64: ऑक्सीजन से आबंधित
कार्बन; 14: ऐल्किल \ce{CH3}।
\end{solution}

\begin{solution}{exo:b2:structure-determination:4}
$(2 \times 8 + 2 - 8)/2 = 5$: एक बेन्ज़ीन वलय (4) और एक \ce{C=O} (1)। उदाहरण के लिए मेथिल बेन्ज़ोएट,
\ce{C6H5COOCH3}, या 4-मेथिलबेन्ज़ोइक अम्ल।
\end{solution}

\begin{solution}{exo:b2:structure-determination:5}
एक असंतृप्ति, एक कीटोन (\qty{1715}{cm^{-1}}, 209 पर रेखा, चतुष्क); एक \ce{CH2} और दो
\ce{CH3}, चार भिन्न कार्बन: ब्यूटेन-2-ओन, \ce{CH3COCH2CH3}।
\end{solution}

\begin{solution}{exo:b2:structure-determination:6}
कार्बन: पेंटेन-2-ओन पाँच रेखाएँ दिखाता है, सममित पेंटेन-3-ओन तीन। प्रोटॉन: पेंटेन-2-ओन में
तीन-प्रोटॉन एकक (\ce{CH3CO}) है; पेंटेन-3-ओन में केवल एक चतुष्क और एक त्रिक। \omterm{def:b2:mass-spec-atomic:spectrum}{द्रव्यमान स्पेक्ट्रम}: पेंटेन-2-ओन
का \omterm{def:b2:mass-spec-atomic:spectrum}{आधार शिखर} 43 (\ce{CH3CO+}) पर है और 58 पर मैकलैफ़र्टी आयन; पेंटेन-3-ओन का \omterm{def:b2:mass-spec-atomic:spectrum}{आधार शिखर} 57
(\ce{C2H5CO+}) पर।
\end{solution}

\begin{solution}{exo:b2:structure-determination:7}
\ce{C4H8O2}, 88: चतुष्क--त्रिक युग्म एक \ce{OCH2CH3} है (4.12 वाला \ce{CH2} ऑक्सीजन पर है),
एकक एक \ce{CH3CO}। एथिल एथेनोएट, \ce{CH3COOCH2CH3}।
\end{solution}

\begin{solution}{exo:b2:structure-determination:8}
1,2-डाइक्लोरोबेन्ज़ीन: तीन रेखाएँ (C1/C2, C3/C6, C4/C5)। 1,4-डाइक्लोरोबेन्ज़ीन: दो (C1/C4, चारों
CH)। DEPT-90 पहले के लिए दो CH रेखाएँ दिखाता है, दूसरे के लिए एक।
\end{solution}

\begin{solution}{exo:b2:structure-determination:9}
C2 एक त्रिविमजनक केंद्र है। C3 के एक या दूसरे प्रोटॉन को ड्यूटेरियम से प्रतिस्थापित करने पर दो
अप्रतिबिंबी त्रिविम समावयवी मिलते हैं, प्रतिबिंबरूप नहीं: दोनों प्रोटॉन अप्रतिबिंबी-स्थानी हैं, भिन्न परिवेशों में।
उनके विस्थापन भिन्न हो सकते हैं और वे आपस में तथा अपने पड़ोसियों से युग्मित होते हैं, अतः
\ce{CH2} साधारण पंचक के बजाय एक जटिल बहुक देता है।
\end{solution}

\begin{solution}{exo:b2:structure-determination:10}
पाँच असंतृप्तियाँ; एक एकल-प्रतिस्थापित वलय (तीन CH रेखाएँ और 137.0 पर एक C रेखा, पाँच \omterm{def:b2:huckel:aromatic}{ऐरोमैटिक}
H); एक कीटोन (200.8, C); एक एथिल समूह (\ce{CH2} 31.8, 2.98 पर चतुष्क, कार्बोनिल के बगल में, \ce{CH3}
8.3)। 1-फ़ेनिलप्रोपेन-1-ओन, \ce{C6H5COCH2CH3}। वलय के साथ संयुग्मन \ce{C=O} के
$\pi$ इलेक्ट्रॉनों को विस्थानीकृत करके आबंध को दुर्बल करता है: पट्ट संतृप्त कीटोन के \qty{1715}{cm^{-1}} से नीचे आ जाता है।
% ledger: ir:CO-ketone
\end{solution}

\begin{solution}{exo:b2:structure-determination:11}
पेंटेनोइक अम्ल: 2500 से \qty{3300}{cm^{-1}} तक अत्यंत चौड़ा \ce{O-H} पट्ट और 11--12
ppm के निकट एक प्रोटॉन। मेथिल ब्यूटेनोएट: ऑक्सीजन से आबंधित कार्बन पर स्थित प्रोटॉनों के परास (3.3--4.5) में
तीन-प्रोटॉन एकक। एथिल प्रोपेनोएट: दो चतुष्क, एक ऑक्सीजन पर और एक कार्बोनिल के बगल में। प्रोपिल
एथेनोएट: कार्बोनिल के बगल में तीन-प्रोटॉन एकक (2.0--2.4) और ऑक्सीजन पर एक त्रिक।
% ledger: ir:OH-acid, nmrr:acid, nmrr:alcohol, nmrr:ketone
\end{solution}

\begin{solution}{exo:b2:structure-determination:12}
दोनों रेखाएँ \omterm{def:b2:structure-determination:dept}{DEPT} में लुप्त हो जाती हैं, अतः दोनों चतुष्क हैं, और केवल \omterm{def:b2:structure-determination:dept}{DEPT} उन्हें अलग नहीं कर सकता। विस्थापन
कर सकते हैं: इस अध्याय के चार्ट में एस्टर कार्बोनिल 167--171 के निकट और \omterm{def:b2:huckel:aromatic}{ऐरोमैटिक} कार्बन 128 और 137 के बीच होते हैं।
166.5 पर वलय कार्बन और 130.6 पर कार्बोनिल दोनों अपने वर्गों से बहुत बाहर होते;
सही निर्धारण इसका उलटा है।
\end{solution}

\begin{solution}{pb:b2:structure-determination:1}
\textbf{1.} $3.1/31.7 \approx 9.8~\%$; $9.8/1.11 \approx 9$ कार्बन।
\textbf{2.} सम द्रव्यमान: कोई नाइट्रोजन नहीं (या दो)।
\textbf{3.} $150 - 108 = 42$: \ce{H10O2} ठीक बैठता है (10 + 32): \ce{C9H10O2}।
\textbf{4.} दस कार्बन $M+1 \approx 11~\%$ देते, $M$ का, 9.8~\% नहीं; और IR एक
कार्बोनिल दिखाता है, जिसके लिए \ce{C-O} पट्ट के अतिरिक्त दूसरा ऑक्सीजन चाहिए।
\textbf{5.} $(2 \times 9 + 2 - 10)/2 = 5$।
\textbf{6.} $108 + 10 \times 1.007825 + 2 \times 15.994915 \approx 150.0681$।
\textbf{7.} संतृप्त एस्टर की स्थिति पर (\qty{1735}{cm^{-1}} के निकट) एक \ce{C=O} तनन।
\textbf{8.} एस्टर का \ce{C-O} तनन।
\textbf{9.} \omterm{def:b2:huckel:aromatic}{ऐरोमैटिक} (या ऐल्कीन) \ce{C-H} आबंध।
\textbf{10.} एक \ce{O-H} समूह: कोई ऐल्कोहॉल नहीं, कोई कार्बोक्सिलिक अम्ल नहीं।
\textbf{11.} नहीं: वलय से संयुग्मित कार्बोनिल निम्नतर तरंग संख्या पर अवशोषण करता है; यहाँ वह वहाँ है जहाँ
संतृप्त एस्टर होते हैं, अतः कोई असंयुग्मक समूह उसे वलय से अलग करता है।
\textbf{12.} 170.70: एस्टर \ce{C=O}; 136.14: प्रतिस्थापी धारण करने वाला वलय कार्बन; 128.56 और 128.24:
\omterm{def:b2:huckel:aromatic}{ऐरोमैटिक} CH; 66.24: ऑक्सीजन से आबंधित \ce{CH2}; 20.82: \ce{CH3}।
\textbf{13.} 66.24, अर्थात् \ce{CH2}।
\textbf{14.} 170.70 और 136.14, हाइड्रोजन-रहित दोनों कार्बन।
\textbf{15.} 128.56 और 128.24।
\textbf{16.} चार: प्रतिस्थापी धारण करने वाला कार्बन, दोनों ऑर्थो, दोनों मेटा और पैरा CH। तीन
CH प्रकार अपेक्षित हैं; उनमें से दो के विस्थापन इस क्षेत्र-सामर्थ्य पर विभेदित होने के लिए बहुत निकट हैं, और वे एक रेखा साझा करते हैं।
\textbf{17.} आवश्यक नहीं: ऊँचाइयाँ गिनतियाँ नहीं हैं (\cref{prop:b2:structure-determination:intensities}),
यद्यपि यहाँ रेखा में दो अतिव्यापी CH प्रकार हो सकते हैं।
\textbf{18.} 7.33: पाँचों \omterm{def:b2:huckel:aromatic}{ऐरोमैटिक} प्रोटॉन; 5.09: \ce{OCH2}; 2.06: \ce{CH3CO}।
\textbf{19.} उनके पड़ोसी परमाणुओं (\ce{CH2} के लिए वलय कार्बन और ऑक्सीजन, \ce{CH3} के लिए कार्बोनिल
कार्बन) पर कोई हाइड्रोजन नहीं है।
\textbf{20.} यह एस्टर के विद्युत-ऋणात्मक ऑक्सीजन से और वलय से आबंधित है: दोनों उसका विपरिरक्षण करते हैं।
\textbf{21.} \ce{C6H5CH2OCOCH3}, बेन्ज़िल एथेनोएट।
\textbf{22.} समावयवी में \ce{CH2} ऑक्सीजन से आबंधित नहीं है, अतः वह 5.09 पर नहीं आ सकता, और
ऑक्सीजन पर स्थित मेथिल 3.3--4.5 परास में आता, कार्बोनिल के बगल में 2.06 पर नहीं।
% ledger: nmrr:alcohol, nmrr:ketone
\textbf{23.} 108: पुनर्विन्यास द्वारा 42 (कीटीन, \ce{CH2=C=O}) की हानि, \ce{C7H8O} का
मूलक धनायन; 91: \ce{C7H7+}, बेन्ज़िल धनायन; 43: \ce{CH3CO+}।
\textbf{24.} कार्बन के सात प्रकार (\ce{C=O}, प्रतिस्थापित वलय कार्बन, ऑर्थो, मेटा और पैरा CH,
\ce{CH2}, \ce{CH3}): \textbf{सात रेखाएँ, जिनमें से एक (\ce{CH2}) नीचे इंगित करती है} DEPT-135 में।
\end{solution}
